% Part: incompleteness
% Chapter: theories-computability
% Section: intepretability

\documentclass[../../../include/open-logic-section]{subfiles}

\begin{document}

\olfileid{inc}{tcp}{itp}

\olsection{$\Th{Q}$ 可在其中被解释的理论是不可判定的}

我们还可以进一步强化这些结果。通俗地说，将语言 $\Lang{L_1}$ \emph{解释}到另一种语言 $\Lang{L_2}$ 中，是指用 $\Lang{L_2}$ 中的公式来定义 $\Lang{L_1}$ 的论域、关系符号和函数符号。尽管我们不在此花篇幅详述，但这一定义是可以严格给出的。

\begin{thm}
  假设 $\Th{T}$ 是某语言中的理论，该语言可以解释算术语言，且 $\Th{T}$ 与 $\Th{Q}$ 的解释是一致的。那么 $\Th{T}$ 是不可判定的。如果 $\Th{T}$ 能证明 $\Th{Q}$ 的公理的解释，那么 $\Th{T}$ 的任何一致扩张都不可判定。
\end{thm}

证明只需对上一条定理的证明稍作修改；可以通过构造反例将 $\Th{Q}$ 与 $\Th{\bar Q}$ 分离开来。可以取 $\Th{ZFC}$，即带选择公理的 Zermelo--Fraenkel（弗兰克尔）集合论，作为一个足够强大、足以展开大量日常数学的公理化基础。在 $\Th{ZFC}$ 中可以定义自然数，且通过这一解释，$\Th{Q}$ 的公理为真。因此我们有

\begin{cor}
$\Th{ZFC}$ 没有可判定的扩张。
\end{cor}

\begin{cor}
$\Th{ZFC}$ 没有完备、一致且可计算公理化的扩张。
\end{cor}

$\Th{ZFC}$ 的语言只有一个二元关系符号 $\in$。（事实上，甚至不需要等词。）所以我们有

\begin{cor}
任何带有一个二元关系符号的语言的一阶逻辑都是不可判定的。
\end{cor}

这一结果可以推广到任何带有两个一元函数符号的语言，因为可用它们来模拟一个二元关系符号。上述结果是最佳的：事实上，仅带有\emph{一元}关系符号和至多一个\emph{一元}函数符号的语言，其一阶逻辑是可判定的。

再补充一点旁注。我们知道，语言 $\Obj 0$, $'$, $+$, $\times$, $<$ 中在标准模型下为真的语句集是不可判定的。事实上，可以用其他符号定义 $<$，进而用 $\times$ 和 $'$ 定义 $+$。因此，语言 $\Obj 0$, $'$, $\times$ 中的真语句集是不可判定的。另一方面，Presburger（普雷斯伯格）已经证明，语言 $\Obj 0$, $'$, $+$ 中在算术语言里为真的语句集是可判定的。然而，该判定过程在计算上是不可行的。

\end{document}